% ==============================================================================
% dxh_closing.tex — DXH workshop: wrap-up
% DXH講座: まとめ
% ==============================================================================

% ------------------------------------------------------------------
\begin{frame}{授業導入に向けて}
  \begin{block}{今日の体験を授業に}
    \begin{itemize}
      \item クラスセット（複数台）での運用イメージ --- 1人1台 / グループ交代のどちらにも対応
      \item 段階的な教材構成 --- 操縦体験 $\to$ 書き込み体験 $\to$ プログラム書き換え、と\\
            難易度を上げながら組み合わせられる
      \item 教材・手順書は今後順次公開予定。本講座の資料もその一部です
    \end{itemize}
  \end{block}
\end{frame}

% ------------------------------------------------------------------
\begin{frame}{アンケートのお願い / リンク集}
  \begin{columns}[c]
    \begin{column}{0.42\textwidth}
      \centering
      \begin{tikzpicture}
        \node[
          rectangle, draw=sfgray, dashed, line width=1pt,
          minimum width=3cm, minimum height=3cm,
          align=center, font=\footnotesize\color{sfgray},
        ] {QR コード\\（当日配布）};
      \end{tikzpicture}
      \vspace{0.3em}

      {\footnotesize アンケートへのご協力をお願いします}
    \end{column}
    \begin{column}{0.55\textwidth}
      \begin{block}{GitHub リポジトリ}
        {\footnotesize\url{https://github.com/M5Fly-kanazawa/stampfly_ecosystem}}
      \end{block}

      \vspace{0.3em}

      \begin{block}{Web フラッシャ}
        {\footnotesize\url{https://m5fly-kanazawa.github.io/stampfly_ecosystem/flash/}}
      \end{block}
    \end{column}
  \end{columns}

  \vspace{0.5em}

  \begin{center}
    \Large\bfseries\textcolor{sfblue}{本日はありがとうございました}
  \end{center}
\end{frame}
